\section{Classical Methods of Pseudorandom Number Generation}

Classical pseudorandom number generators are based on simple mathematical formulas, typically involving modular arithmetic and recurrence relations. Despite their simplicity, they form the foundation for many modern generators.

In this section, we present three important types of classical generators: the Linear Congruential Generator (LCG), the Lehmer generator, and combined generators.

\subsection{Linear Congruential Generator (LCG)}

The Linear Congruential Generator is one of the oldest and simplest PRNGs. It generates a sequence of numbers using the recurrence relation:

$$x_{n+1} = (a x_n + c) \mod m$$

where:
\begin{itemize}
    \item $x_n$ is the current state,
    \item $a$ is the multiplier,
    \item $c$ is the increment,
    \item $m$ is the modulus,
    \item $x_0$ is the seed.
\end{itemize}

The sequence $(x_n)$ takes values in the set $\{0, 1, 2, ..., m-1\}$.

\subsubsection{Period of LCG}

The maximum possible period of an LCG is $m$. This is achieved if and only if the following conditions are satisfied:

\begin{itemize}
    \item $\gcd(c, m) = 1$,
    \item $a - 1$ is divisible by all prime factors of $m$,
    \item if $m$ is divisible by 4, then $a - 1$ is also divisible by 4.
\end{itemize}

These conditions are known as the Hull-Dobell Theorem.

\subsubsection{Properties and Limitations}

LCGs are easy to implement and computationally efficient. However, they suffer from several drawbacks:

\begin{itemize}
    \item relatively short periods for poorly chosen parameters,
    \item visible patterns in higher dimensions,
    \item limited statistical quality compared to modern generators.
\end{itemize}

\subsubsection{Example of LCG}

Example with parameters:
$m = 9$, $a = 2$, $c = 1$, $x_0 = 1$

\begin{table}[h]
\centering
\small
\begin{tabular}{|c|c|}
\hline
$n$ & $x_n$ \\
\hline
0 & 1 \\
1 & 3 \\
2 & 7 \\
3 & 6 \\
4 & 4 \\
5 & 0 \\
6 & 1 \\
\hline
\end{tabular}
\caption{Example sequence generated by LCG}
\end{table}

\subsection{Lehmer Generator}

The Lehmer generator is a special case of the LCG where the increment $c = 0$. Its recurrence relation is:

$$x_{n+1} = (a x_n) \mod m$$

It is also known as the multiplicative congruential generator.

\subsubsection{Properties}

The Lehmer generator has a simpler structure than the general LCG, but requires careful choice of parameters.

To achieve a long period, the modulus $m$ is often chosen as a prime number, and the multiplier $a$ should be a primitive root modulo $m$.

In that case, the generator achieves the maximum period:

$$T = m - 1$$

\subsubsection{Advantages and Disadvantages}

\begin{itemize}
    \item simpler structure than LCG,
    \item good statistical properties for well-chosen parameters,
    \item still vulnerable to patterns and predictability.
\end{itemize}

\subsection{Combined Generators}

Combined generators are constructed by combining two or more simpler generators (such as LCGs) in order to improve statistical properties and increase the period.

A typical combined generator may take the form:

$$x_n = (x_n^{(1)} - x_n^{(2)}) \mod m$$

where $x_n^{(1)}$ and $x_n^{(2)}$ are sequences generated by two independent generators.

\subsubsection{Period}

If the individual generators have periods $T_1$ and $T_2$, the combined generator can achieve a period close to:

$$T \approx \mathrm{lcm}(T_1, T_2)$$

where $\mathrm{lcm}$ denotes the least common multiple.

\subsubsection{Advantages}

\begin{itemize}
    \item significantly longer period,
    \item improved statistical properties,
    \item reduced correlation between generated values.
\end{itemize}

\subsection{Summary}

Classical generators such as LCG and Lehmer are simple and efficient but have limitations in terms of statistical quality. Combined generators address some of these issues by increasing the period and improving randomness properties.

These methods form the basis for understanding more advanced pseudorandom number generators.